% Part: normal-modal-logic
% Chapter: syntax-and-semantics
% Section: introduction

\documentclass[../../../include/open-logic-section]{subfiles}

\begin{document}

\olfileid[id]{nml}{syn}{int}

\olsection{Pengantar}

Logika modal membahas \emph{proposisi modal} dan relasi perikutan di antara
proposisi-proposisi tersebut. Berikut ini beberapa contoh proposisi modal:
\begin{enumerate}
\item Niscaya bahwa $2+2=4$.
\item Niscaya mungkin bahwa besok akan turun hujan.
\item Jika niscaya mungkin bahwa~$!A$, maka mungkin bahwa~$!A$.
\end{enumerate}
Kemungkinan dan keniscayaan bukanlah satu-satunya modalitas: perangkai uner
lainnya juga digolongkan sebagai modalitas, misalnya, ``seharusnya demikian
bahwa~$!A$,'' ``akan demikian bahwa~$!A$,'' ``Dana mengetahui bahwa~$!A$,''
atau ``Dana meyakini bahwa~$!A$.''

Logika modal pertama kali muncul dalam \emph{De Interpretatione} karya
Aristoteles: dialah yang pertama kali mengamati bahwa keniscayaan
mengimplikasikan kemungkinan, tetapi tidak sebaliknya; bahwa kemungkinan dan
keniscayaan dapat saling didefinisikan; bahwa jika $!A \land !B$ mungkin
benar, maka $!A$ mungkin benar dan $!B$ mungkin benar, tetapi tidak
sebaliknya; dan bahwa jika $!A \to !B$ niscaya, maka apabila $!A$ niscaya,
$!B$ juga niscaya.

Pendekatan modern pertama terhadap logika modal adalah karya C.~I. Lewis,
yang berpuncak pada buku Lewis dan Langford, \emph{Symbolic Logic}
(1932). Lewis \& Langford tidak puas dengan representasi implikasi melalui
kondisional material: $!A \lif !B$ merupakan pengganti yang buruk bagi
``$!A$ mengimplikasikan $!B$.'' Sebagai gantinya, mereka mengusulkan agar
implikasi dicirikan sebagai ``Niscaya, jika $!A$ maka $!B$,'' yang
disimbolkan sebagai $!A \strictif !B$. Dalam upaya memilah berbagai
sifatnya, Lewis mengidentifikasi lima sistem modal yang berbeda, \Log{S1},
\ldots, \Log{S4}, \Log{S5}; dua sistem terakhir masih digunakan hingga kini.

Pendekatan Lewis dan Langford sepenuhnya bersifat \emph{sintaktis}: mereka
mengidentifikasi aksioma dan aturan yang masuk akal serta menyelidiki apa yang
dapat dibuktikan dengan sarana tersebut. Pendekatan semantik lama tak
terjangkau, sampai Rudolf Carnap melakukan upaya pertama dalam \emph{Meaning
  and Necessity} (1947) dengan menggunakan gagasan \emph{deskripsi keadaan},
yakni kumpulan kalimat atomik (kalimat yang ``benar'' dalam deskripsi keadaan
tersebut). Setelah memperluas definisi kebenaran ke kalimat sebarang $!A$,
Carnap mendefinisikan $!A$ sebagai \emph{benar secara niscaya} jika kalimat
itu benar dalam semua deskripsi keadaan. Pendekatan Carnap tidak dapat
menangani modalitas \emph{berulang}, sebab kalimat berbentuk ``Mungkin secara
niscaya \ldots mungkin $!A$'' selalu tereduksi ke modalitas terdalam.

Terobosan besar dalam semantik modal hadir melalui artikel Saul Kripke ``A
Completeness Theorem in Modal Logic'' (JSL 1959). Kripke mendasarkan
pekerjaannya pada gagasan Leibniz bahwa suatu pernyataan benar secara niscaya
jika pernyataan itu benar ``di semua dunia mungkin.'' Namun, gagasan ini
menghadapi kelemahan yang sama seperti pendekatan Carnap, sebab kebenaran
suatu pernyataan pada dunia $w$ (atau deskripsi keadaan $s$) sama sekali tidak
bergantung pada $w$. Karena itu, Kripke mengandaikan bahwa dunia-dunia
dihubungkan oleh suatu \emph{relasi aksesibilitas} $R$, dan bahwa pernyataan
berbentuk ``Niscaya $!A$'' bernilai benar pada dunia $w$ jika dan hanya jika
$!A$ bernilai benar pada semua dunia $w'$ yang \emph{dapat diakses dari}
$w$. Semantik yang menyediakan suatu versi pendekatan ini disebut semantik
Kripke dan memungkinkan perkembangan logika-logika modal (dalam bentuk
jamak) yang begitu pesat.

Ketika ditafsirkan dengan semantik Kripke, logika modal memperlihatkan kepada
kita bagaimana \emph{struktur relasional} tampak ``dari dalam.'' Struktur
relasional hanyalah suatu himpunan yang dilengkapi dengan relasi biner
(misalnya, himpunan siswa di kelas yang diurutkan menurut nomor jaminan sosial
mereka merupakan struktur relasional). Namun, struktur relasional sebenarnya
hadir dalam segala macam ranah: selain kemungkinan relatif keadaan dunia,
kita dapat memiliki keadaan epistemik suatu agen yang dihubungkan oleh
kemungkinan epistemik, atau keadaan suatu sistem dinamis beserta transisi
keadaannya, dan sebagainya. Logika modal dapat digunakan untuk memodelkan
semuanya: yang pertama memberi kita logika modal aletik biasa; yang lainnya
memberi kita logika epistemik, logika dinamis, dan sebagainya.

Kita berfokus pada satu sudut pandang tertentu yang dikenal oleh para ahli
logika modal sebagai ``teori korespondensi.'' Salah satu penemuan awal Kripke
yang paling penting ialah bahwa banyak sifat relasi aksesibilitas~$R$
(apakah relasi itu transitif, simetris, dan sebagainya) dapat dicirikan
\emph{dalam bahasa modal} itu sendiri melalui ``skema modal'' yang sesuai.
Sebagai contoh, para ahli logika modal mengatakan bahwa refleksivitas $R$
``berkorespondensi'' dengan skema ``Jika niscaya~$!A$, maka~$!A$''. Kita
terutama mengkaji teori korespondensi sejumlah sistem klasik logika modal
(misalnya, \Log{S4} dan \Log{S5}) yang diperoleh melalui kombinasi skema
\Ax{D}, \Ax{T}, \Ax{B}, \Ax{4}, dan~\Ax{5}.

\end{document}
